\chapter{Financial Plan}
\label{chapter:ent-financial}

This chapter presents the financial projections for the venture. Unit economics and operational cost assumptions draw on the supply chain analysis in \Cref{sec:ent-supply-chain} and the adoption scenarios in \Cref{chapter:ent-gtm}. The first four sections cover the core financial model; the funding strategy in \Cref{sec:ent-funding} addresses how the venture is capitalised through to breakeven.


\section{Key Assumptions}
\label{sec:ent-assumptions}

\todo{Present a clear assumptions table covering:
\begin{itemize}
    \item Unit cost and retail price (by volume tier)
    \item Subscription pricing and churn rate
    \item Customer acquisition timeline (from GTM chapter)
    \item Headcount and compensation schedule (from operations)
    \item Regulatory and development cost estimates
    \item Working capital assumptions (inventory, receivables, payables)
\end{itemize}
This table should be the single source of truth for all subsequent projections. The reader should be able to change an assumption and understand how it flows through.}


\section{Profit and Loss Projection}
\label{sec:ent-pnl}

\todo{5-year P\&L with:
\begin{itemize}
    \item Revenue breakdown (hardware, subscriptions, research licences)
    \item Cost of goods sold
    \item Gross profit and margin
    \item Operating expenses by category (R\&D, sales, regulatory, G\&A)
    \item EBITDA
\end{itemize}
Include a table and a chart showing revenue vs costs over time. Explain the path to breakeven.}


\section{Cash Flow Projection}
\label{sec:ent-cash-flow}

\todo{5-year cashflow with:
\begin{itemize}
    \item Operating cashflow (EBITDA adjusted for working capital changes)
    \item Investing cashflow (capex, tooling, software capitalisation)
    \item Financing cashflow (equity injections)
    \item Closing cash balance
\end{itemize}
Highlight the minimum cash point and what it implies for funding requirements.}


\section{Sensitivity Analysis}
\label{sec:ent-sensitivity}

\todo{Stress-test key assumptions:
\begin{itemize}
    \item Base case, conservative, and optimistic scenarios
    \item What happens if adoption is 50\% slower?
    \item What happens if BOM costs are 30\% higher?
    \item What happens if regulatory takes an extra year?
    \item Tornado chart or similar showing which assumptions have the largest impact on breakeven timing
\end{itemize}
This section demonstrates rigour and shows which variables the business is most sensitive to.}


\section{Funding and Investment Strategy}
\label{sec:ent-funding}

\todo{Eesh: Draft this section.}

\subsection{Development Cost Estimate}

\todo{Estimate total development costs across phases.}

\subsection{Funding Requirements}

\todo{Total capital required to reach breakeven, timing and staging of funding rounds.}

\subsection{Funding Sources}

\todo{Evaluate candidate funding sources: grants, angel, VC, non-dilutive.}

\subsection{Investor Proposition}

\todo{Summarise the investment case: return profile, value inflection points, comparable exits.}
